% =======================================================================================================
% Capítulo 1 - Introdução (versão enxuta)
% Convenção de siglas: na primeira ocorrência, nome em português, termo em inglês (quando houver) e
% sigla; depois, só a sigla.
% =======================================================================================================

\section{Problema e motivação}
\label{sec:formulacao}
\label{sec:motivacao}

A cavitação é um dos fenômenos que mais comprometem bombas centrífugas e turbinas hidráulicas. Quando a pressão local cai abaixo da pressão de vapor do líquido, formam-se bolhas que colapsam ao atingir regiões de pressão mais alta, causando erosão, trincas, ruído e perda de desempenho \cite{brennen1995,liu2023review}. Detectá-la cedo, sem interromper a operação, é relevante para a segurança e para a manutenção preditiva \cite{mba2006development}.

A Emissão Acústica (EA) é uma técnica adequada a essa tarefa: ela capta as ondas elásticas geradas pelo colapso das bolhas e pela erosão incipiente e é sensível já nas fases iniciais da falha \cite{Neill1997Detection,Alfayez2005ndt,FernandezOsete2024}. Essa sensibilidade tem um custo. Os sinais de EA ocupam a faixa de dezenas de quilohertz a cerca de 1\,MHz e precisam ser amostrados a vários megahertz \cite{grosse2008acoustic}. No Projeto MOMEA (Monitoramento Online de Rotores de Turbinas pelo Método de Emissão Acústica), do Itaipu Parquetec e da Itaipu Binacional, a aquisição contínua a 8\,MSa/s por canal produziu dezenas de gigabytes por hora e exigiu compactação e gravação contínuas no computador, com a análise feita só depois da gravação (Seção~\ref{sec:mm_antecedentes}).

A literatura de sistemas embarcados aponta uma saída: deslocar para o \textit{hardware} reconfigurável a etapa que reduz o sinal a poucos atributos, de modo que apenas a informação relevante siga adiante \cite{kang2015,wirtz2018,toumi2022fpga}. Os \textit{Field-Programmable Gate Arrays} (FPGAs) combinam paralelismo e latência determinística, o que os torna adequados a essa função \cite{trimberger2015three,Shawahna2019}. Na prática da EA, a redução mais estabelecida é a detecção de \textit{hits}: o sinal só é registrado quando ultrapassa um limiar, e cada evento é resumido por amplitude, duração, contagens e energia \cite{unnthorsson2013,vallen2022}. Resta saber se esses descritores bastam para reconhecer a condição de operação da máquina por Aprendizado de Máquina (AM) e se essa classificação pode ser executada no próprio FPGA.

\section{Objetivos}
\label{sec:objetivos}

O objetivo geral é desenvolver e avaliar uma arquitetura de processamento em FPGA para a detecção de falhas em motobombas por EA, na qual os sinais são reduzidos a registros de \textit{hits} na lógica programável e a condição de operação é inferida por uma Rede Neural Multicamadas (\textit{Multilayer Perceptron}, MLP) a partir desses registros. Os objetivos específicos são:
\begin{enumerate}[label=\textbf{\arabic*.}]
	\item implementar a extração de \textit{hits} em um FPGA Artix-7, com registro compacto por evento e as mesmas convenções temporais de um sistema comercial;
	\item avaliar o extrator em ensaio simultâneo com o sistema comercial Vallen AMSY-6, quanto a detecção, descritores, integridade da comunicação e redução de dados;
	\item treinar MLPs com descritores de \textit{hits} de uma bomba em bancada e avaliar a generalização com validações que separem corridas e sessões independentes;
	\item quantizar o modelo em int8 e implementá-lo no processador MicroBlaze e em lógica dedicada, verificando por simulação a fidelidade, a latência e o uso de recursos;
	\item validar o método em uma base pública de defeitos em rolamentos.
\end{enumerate}

\section{Contribuições e organização do texto}
\label{sec:contribuicoes}
\label{sec:organizacao}

As contribuições são: um extrator de \textit{hits} de EA em FPGA de baixo custo, equivalente a um instrumento comercial na faixa ensaiada e com redução de cerca de 770 vezes no volume de dados; a demonstração de que descritores de \textit{hits} permitem estimar a condição de operação de uma bomba, com uma análise crítica de quando o resultado se generaliza; uma cadeia de inferência int8 com fonte única para MicroBlaze e Vitis HLS, verificada em seis níveis até a simulação RTL do processador; e a validação externa do método na base CWRU.

O Capítulo~\ref{chap:fundamentacao} apresenta os fundamentos de EA e de detecção de \textit{hits}, de redes MLP e de FPGAs. O Capítulo~\ref{chap:metodos} descreve a arquitetura proposta, os ensaios, os dados e os métodos de classificação, de implementação embarcada e de validação. O Capítulo~\ref{chap:resultados} apresenta e discute os resultados, e o Capítulo~\ref{chap:conclusoes} reúne as conclusões e os trabalhos futuros.
